\chapter{User Guide}

This appendix provides a step-by-step guide for setting up, running, and evaluating the synthetic ICU trajectory generation pipeline. All commands assume a Unix-like terminal (macOS or Linux) and are run from the project root directory.

%==============================================================================
\section{Prerequisites}
%==============================================================================

\begin{itemize}\tightlist
    \item \textbf{Python 3.10 or later.}
    \item \textbf{uv} (recommended) or \textbf{pip} for dependency management. Install uv via:
    \begin{verbatim}
    curl -LsSf https://astral.sh/uv/install.sh | sh
    \end{verbatim}
    \item \textbf{MIMIC-III access.} A credentialed PhysioNet account with a signed data use agreement is required to obtain the source data. See \url{https://mimic.physionet.org/gettingstarted/access/}.
    \item \textbf{MIMIC-Extract output.} The pipeline expects hourly ICU tables produced by a MIMIC-Extract style workflow, placed under the project data directory.
\end{itemize}

%==============================================================================
\section{Installation}
%==============================================================================

Clone the repository and install dependencies:

\begin{verbatim}
git clone <repository-url>
cd lexisflow
uv sync            # installs all dependencies from pyproject.toml
\end{verbatim}

Alternatively, with pip:

\begin{verbatim}
pip install -e .
\end{verbatim}

To verify the installation, run the test suite:

\begin{verbatim}
uv run pytest
\end{verbatim}

All unit tests should pass.

%==============================================================================
\section{Configuration}
%==============================================================================

Runtime configuration is now centralised in Python dataset presets plus CLI overrides.
The authoritative defaults are defined in \texttt{src/lexisflow/config/datasets.py}.

\begin{itemize}\tightlist
    \item \textbf{Dataset paths (\texttt{DatasetConfig}):} script directory, processed CSV paths, preprocessor artifact paths, transformed cache location, results CSV, and sweep artifact directory.
    \item \textbf{Sweep defaults (\texttt{SweepDefaults}):} \texttt{nt} grid, \texttt{n\_noise} grid, training row cap, synthetic sample count, and privacy-evaluation row cap.
    \item \textbf{Split policy (\texttt{SplitConfig}):} patient-level test fraction, holdout fraction, and split shuffle seed.
    \item \textbf{Profiles (\texttt{sweep\_profiles}):} each dataset exposes \texttt{full} and \texttt{smoke} presets. Scripts resolve defaults from the selected profile, then apply explicit CLI overrides.
\end{itemize}

Practical rule: treat \texttt{datasets.py} as the single source of truth for default
paths, sweep budgets, and split ratios.

%==============================================================================
\section{Pipeline Walkthrough}
%==============================================================================

The pipeline is now executed through a single command that orchestrates data
preparation, preprocessor fitting, and the sweep run.

\subsection{Single Command Pipeline}

\begin{verbatim}
uv run python scripts/run_sweep.py --dataset mimic
uv run python scripts/run_sweep.py --dataset challenge2012
uv run python scripts/run_sweep.py --dataset mimic --reset
uv run python scripts/run_sweep.py --dataset mimic --profile smoke
\end{verbatim}

\textbf{Outputs (MIMIC):}
\begin{itemize}\tightlist
    \item \texttt{data/processed/hour0\_data.csv}
    \item \texttt{data/processed/autoregressive\_data.csv}
    \item \texttt{artifacts/hour0\_preprocessor.pkl}
    \item \texttt{artifacts/preprocessor\_full.pkl}
    \item \texttt{results/sweep\_results.csv}
\end{itemize}

\textbf{Outputs (Challenge 2012):}
\begin{itemize}\tightlist
    \item \texttt{data/challenge2012/processed/hour0\_data.csv}
    \item \texttt{data/challenge2012/processed/autoregressive\_data.csv}
    \item \texttt{artifacts/challenge2012/hour0\_preprocessor.pkl}
    \item \texttt{artifacts/challenge2012/autoregressive\_preprocessor.pkl}
    \item \texttt{results/challenge2012\_sweep\_results.csv}
\end{itemize}

\subsection{Step 4: Evaluate Quality}

Quality metrics are computed automatically for every sweep cell (KS statistics, correlation Frobenius norm, clinical range violations, privacy diagnostics). Use \texttt{scripts/run\_sweep.py --dataset mimic} (or \texttt{--dataset challenge2012}) for end-to-end execution.

%==============================================================================
\section{Hyperparameter Sweep}
%==============================================================================

For systematic evaluation across configurations, use the unified sweep entrypoint:

\begin{verbatim}
uv run python scripts/run_sweep.py --dataset mimic
uv run python scripts/run_sweep.py --dataset challenge2012
uv run python scripts/run_sweep.py --dataset mimic --reset
\end{verbatim}

Key options:
\begin{itemize}\tightlist
    \item \texttt{--dataset}: choose \texttt{mimic} or \texttt{challenge2012}.
    \item \texttt{--profile}: choose \texttt{full} or \texttt{smoke}; forwarded to the dataset-specific sweep driver.
    \item \texttt{--reset}: clear generated data, preprocessors, cache and sweep outputs before rerun.
\end{itemize}

The sweep trains both hour-0 and autoregressive models for each
$(n_t, n_{\text{noise}})$ configuration, generates synthetic samples, runs
multi-task TSTR and quality/privacy metrics, and appends to the dataset-specific
results CSV (\texttt{results/sweep\_results.csv} for MIMIC,
\texttt{results/challenge2012\_sweep\_results.csv} for Challenge 2012). The
sweep is resumable: existing rows are skipped on rerun.

Dataset sweep drivers (\texttt{scripts/mimic/run\_sweep.py} and
\texttt{scripts/challenge2012/run\_sweep.py}) also support direct execution with
fine-grained overrides:
\begin{itemize}\tightlist
    \item \texttt{--profile}, \texttt{--train-rows}, \texttt{--synth-samples}
    \item \texttt{--nt-values}, \texttt{--noise-values}, \texttt{--n-jobs}
    \item \texttt{--model-type} (\texttt{hs3f}, \texttt{forest-flow}, \texttt{ctgan})
    \item \texttt{--refresh-transformed-cache}, \texttt{--skip-privacy-metrics}, \texttt{--privacy-max-rows}
\end{itemize}

To visualise sweep results:

\begin{verbatim}
uv run python scripts/common/analyze_sweep.py
uv run python scripts/common/analyze_sweep.py --dataset challenge2012
uv run python scripts/common/analyze_sweep.py --results-path results/challenge2012_sweep_results.csv
\end{verbatim}

\texttt{analyze\_sweep.py} options:
\begin{itemize}\tightlist
    \item \texttt{--dataset}: selects default results path from dataset config.
    \item \texttt{--results-path}: explicit input CSV override.
    \item \texttt{--output-dir}: explicit plot output directory override.
\end{itemize}

%==============================================================================
\section{Running Tests}
%==============================================================================

The project includes 54 unit tests covering quality metrics, privacy metrics, TSTR framework, autoregressive transformation, feature utilities, model training, and sampling:

\begin{verbatim}
uv run pytest                    # run all tests
uv run pytest -v                 # verbose output
uv run pytest src/lexisflow/models/   # run model tests only
\end{verbatim}
